d) Of course, $\mathfrak g$ denotes the Lie algebra of $G$, and $T\to\GL(\mathfrak g)$ the homomorphism induced by the adjoint representation of $G$. One trivially has
\[
T\cap\Centr_G\subset\Ker(T\to\GL(\mathfrak g));
\]
the proof of the reverse inclusion is the same as in 4.7 d), and we do not repeat it here. Let $Z$ be the group in question; as $\Ker(T\to\GL(\mathfrak g))$ it is of multiplicative type (\cf for example IX~6.8). To prove that it is a reductive center of $G$, one is reduced by b) to the case where $S$ is the spectrum of a field. Let then $Z'$ be the reductive center (which exists by a)); one obviously has $Z'\subset T\cap\Centr(G)=Z$; on the other hand, since $Z$ is a central subgroup of multiplicative type contained in the smooth connected affine subgroup $T$, it follows from a) that $Z'\supset Z$, hence $Z'=Z$. \hfill Q.E.D.

e) Under the conditions of 7.1 f), put $Z=\Ker u$ and suppose that for every algebraically closed field $k$ over $S$, $Z_k$ is contained in a maximal torus of $G_k$. Then one observes easily that the map $T'\mto u^{-1}(T')$ induces a one-to-one correspondence between the set of maximal tori of $G'$ and the set of maximal tori of $G$. Applying this to situation 8.8 e), the desired conclusion follows at once.
\end{proof}
\begin{remarks}{8.9}\label{XII.8.9}
a) The proof given of 8.8 is independent of the results of \No4, and in particular that of 8.8 a) does not use 4.4 (whose proof is somewhat laborious).

b) One sees easily that the subgroup $Z$ of $G$ considered in 8.6 is always central (whether representable or not), and it may be tempting to call it the reductive center of $G$ in all cases.

c) One can also generalize 4.9; one finds the following statement: let $G$ be a smooth connected algebraic group over an algebraically closed field; for $G$ to be an extension of an abelian variety by a smooth connected unipotent algebraic group (\ie for the reductive rank of $G$ to be zero, \cf 8.3), it is necessary and sufficient that the reductive center of $G$ be reduced to the identity group, and that the Lie algebra of $G$ be nilpotent.
\end{remarks}

\sgasection{9}{Supplement: action of a group scheme and fixed points}
The aim of this section, added in January 2008, is to study the fixed points under the action of a group scheme.

\numpara{9.1}\textbf{Representability of the fixed-point functor.}---Let $S$ be a scheme, $G$ an $S$-group scheme acting on an $S$-scheme $X$. One defines the subfunctor $X^G$ of $X$ of fixed points of $X$ under $G$: for every $S$-scheme $T$, $X^G(T)$ is the subset of $X(T)$ formed by the fixed points under $G$ (VIII.6.e).

Let us recall the notion of an ``$S$-pure scheme'' introduced in ([G-R], §~3.3, [R]). Let $G$ be a flat $S$-group scheme of finite presentation. Since the fibers of $G$ have no embedded components, the notion of purity for $G$ takes the following form. The scheme $G$ is $S$-pure if for every strictly henselian local $S$-scheme $T$, with closed point $t$, every generic point $\widetilde x$ of a fiber of $G\times_S T$ specializes to a point of $G_t$, \ie the closure of $\widetilde x$ in $G\times_S T$ meets $G_t$.
% typo: stray quote preceding the source's opening quotation mark -> omitted.

\par\medskip\noindent\textbf{Examples.}---
\begin{enumerate}[label=(\arabic*)]
\item If $G$ is quasi-finite and separated over $S$, $G$ is $S$-pure if and only if $G$ is finite over $S$.
\item If $S=\Spec(R)$ and $G=\Spec(A)$, $G$ is $S$-pure if and only if $A$ is a projective $R$-module ([G-R], 3.3.5). In particular $G$ is $S$-pure if $G$ is diagonalizable.
\item A group scheme of multiplicative type is $S$-pure, for the notion of purity is local for the étale topology on $S$.
\item The scheme $G$ is $S$-pure if $G$ is $S$-proper or if the fibers of $G$ are irreducible or if $S$ is semilocal artinian.
\item In particular, a reductive $S$-group scheme $G$ is $S$-pure.
\end{enumerate}
\begin{proposition}{9.2}\label{XII.9.2}
Let $G$ be an $S$-group scheme of finite presentation, $S$-flat and $S$-pure, acting on an $S$-scheme $X$ of finite presentation. Then the functor $X^G$ of fixed points of $X$ under $G$ is representable by a closed subscheme of $X$, of finite presentation over $S$.
\end{proposition}
Indeed, let $a$ be in $X(S)$ and let $H$ be the group subscheme of $G$ that is the stabilizer of $a$. Then $a$ is fixed under $G$ if and only if $H=G$. Now the subfunctor $C$ of $S$ of coincidences (more precisely, if $T$ is an $S$-scheme, $C(T)=\{\varnothing\}$ if $H\times_S T\to G\times_S T$ is bijective, and $C(T)=\varnothing$ otherwise) of $H$ with $G$ is representable by a closed subscheme of $S$, defined by a sheaf of ideals of finite type ([G-R], 4.1.1). One applies this result with $S=X$, taking for $a$ the universal point of $X$ (\ie the identity in $X(X)=\Hom_S(X,X)$).

\numpara{9.3}\textbf{Infinitesimal obstruction.}---Under the hypothesis that $G$ is a flat $S$-group scheme of finite presentation acting (on the left) on a smooth $S$-prescheme $X$, we shall now study the formal smoothness of the functor $X^G$.

One supposes here that $S$ is endowed with a closed subscheme $S_0$ defined by a sheaf of quasi-coherent ideals $\mathscr I$ of square zero. One writes $j:S_0\to S$, $G_0=G\times_S S_0$, $X_0=X\times_S S_0$.
% typo: "un faisceaux" -> "un faisceau".

Let $\epsilon_0$ be a point of $X^G(S_0)$ lifting to a point of $X(S)$. We shall study the obstruction to lifting $\epsilon_0$ to a point of $X^G(S)$. Since $X$ is smooth, one knows that the liftings of $\epsilon_0$ to $X(S)$ form a trivial principal homogeneous space under the abelian group
\[
\Hom_{\OO_{S_0}}(\epsilon_0^*(\Omega^1_{X_0/S_0}),\mathscr I)
\]
(III.0.2, 0.3). One then puts
\[
L_0=\uHom_{\OO_{S_0}}(\epsilon_0^*(\Omega^1_{X_0/S_0}),\mathscr I),
\]
which is an $\OO_{S_0}$-module. Moreover, since $\epsilon_0$ is fixed under $G$, $L_0$ is naturally a $G_0$-$\OO_{S_0}$-module. One denotes by $\rho_0:G_0\to\underline{\Aut}_{S_0}(L_0)$ (resp. $\rho:G\to\underline{\Aut}_S(j_*L_0)$) the associated representation.
% typo: "pour le sous le groupe" -> "sous le groupe".
\begin{lemma}{9.4}\label{XII.9.4}
There exists a certain class $c(\epsilon_0)\in H^1(G,j_*L_0)\simeq H^1(G_0,L_0)$, defined canonically by $\epsilon_0$, such that $\epsilon_0$ lifts to $X^G(S)$ if and only if $c(\epsilon_0)=0$.
\end{lemma}
The adjunction $H^1(G,j_*L_0)\simeq H^1(G_0,L_0)$ is that of lemma III.1.1.2. We shall work with the small fppf site over $S$. For every scheme $T$ flat and of finite presentation over $S$, write $G_T=G\times_S T$, $X_T=X\times_S T$ for the corresponding objects over $T$. Consider the sheaf $A$ on the small fppf site of $S$ such that for every $T$ flat and of finite presentation over $S$, one has:
\[
A(T)=\left\{\text{set of liftings of }\epsilon_0\times_S T\text{ to }X(T)\right\}.
\]
Since the formation of $L_0$ commutes with flat base changes, the preceding considerations indicate that $A(T)$ is a trivial principal homogeneous space under the abelian group $H^0(T,j_*L_0)$. Again because $\epsilon_0$ is fixed, for every $T$ flat and of finite presentation over $S$, and every $g$ in $G(T)$, $g$ acts by affine automorphisms on $A(T)$, compatibly with the action of $G$ on $j_*L_0$. That is:
\[
g(a_T+v)=g(a_T)+\rho(g)(v),\qquad v\in H^1(T,j_*L_0).
\]
Since $G$ is flat of finite type, one can apply these considerations by taking $T=G$ and for $g$ the universal point $\mathrm{id}_G$ of $G$. One thus obtains an action of the fppf sheaf $G$ on $A$.
% typo?: H^1 rather than H^0 in the affine-action formula retained.

Let $a$ be a lifting of $\epsilon_0$ to $X(S)$. Denote by $g^\sharp$ the universal point of $G$ and define $v^\sharp\in H^0(G,j_*L_0)$ by
\[
\rho(g^\sharp)(v^\sharp)=g^\sharp\cdot a_G-a_G.
\]
For every $S$-prescheme $Y$, put $z(Y):G(Y)\to H^0(Y,j_*L_0)$, $g\mto v^\sharp(g)$. This defines the $1$-cocycle $z$ in $Z^1(G,L)$ (exposé I, §~5).
% typo?: L in Z^1(G,L) not specified, retained.

Its class $c(\epsilon_0)$ in $H^1(G,j_*L_0)$ does not depend on the choice of $a$. In particular, if $\epsilon_0$ lifts to $X^G(S)$, one has $c(\epsilon_0)=0$. Conversely, if $c(\epsilon_0)=0$, then there exists $w\in H^0(S,j_*L_0)$ such that $z(Y)(g)=g\cdot w_Y-w_Y$ for every $S$-prescheme $Y$ and every $g\in G(Y)$. Applying this to $Y=G$ and to $g^\sharp$, one concludes that $a-w\in X^G(S)$. We have therefore established that $\epsilon_0$ lifts to $X^G(S)$ if and only if $c(\epsilon_0)=0$.
% typo: missing "pour" before "tout S-préschéma" -> supplied in English.
\begin{remark}{9.5}\label{XII.9.5}
In the case where $G$ is affine and flat of finite type over affine $S$, one can refine this obstruction to a class $\widetilde c(\epsilon_0)\in H^1_{G_0}(X_0,L_0)$, where $H^1_{G_0}(X_0,L_0)$ is the $G_0$-equivariant cohomology group defined by Wevers ([W], app.~C). In this case, it is not necessary to suppose that $\epsilon_0$ lifts to $X(S)$.
% typo?: name Wevers in prose differs from Wewers in bibliography; retained.
\end{remark}

\numpara{9.6}\textbf{Smoothness of fixed points.}---
\begin{theorem}{9.7}\label{XII.9.7}
Let $G$ be a flat $S$-group scheme of finite type over a noetherian base $S$, acting on a smooth $S$-scheme $X$. Suppose that $G$ is $S$-pure, and for every geometric point $s$ over $S$ (with algebraically closed residue field), one has
\[
H^1\bigl(G_s,(\Omega^1_{X_s})^*\bigr)=0.
\]
Then the functor $X^G$ of fixed points of $X$ under $G$ is representable by a closed subscheme of $X$ smooth over $S$.
\end{theorem}
Let us prove theorem 9.7. One knows that $X^G$ is representable by a closed subscheme of $X$ of finite presentation over $S$. One can suppose $S=\Spec(A)$ local. Following the smoothness criterion (SGA~1 III.3.1.iii.bis), it suffices to verify the formal smoothness of $X^G$ for $S'=\Spec(A')$ and the closed subscheme $S'_0=\Spec(A'/\mathscr I')=\Spec(A'_0)$, where $A'$ is an artinian local ring and $\mathscr I'$ a square-zero ideal of $A'$. By reduction, one is reduced to the case where $\mathscr I'$ is annihilated by the maximal ideal of $A'$. In particular, if $k$ denotes the residue field of $A'$, $\mathscr I'$ is a $k$-vector space.

One therefore gives oneself an element $\epsilon'_0\in X^G(S'_0)$ which one wishes to lift to $X^G(S')$. One then writes $X'=X\times_S S'$, $G'=G\times_S S'$, $X'_0=X\times_S S'_0$, $G'_0=G\times_S S'_0$. Since $S'$ is affine and $X'$ is smooth over $S'$, $\epsilon'_0\in X^G(S'_0)$ lifts to $X(S')$. According to lemma 9.4, the class
\[
c(\epsilon_0)\in H^1(G'_0,L'_0)
\]
is the obstruction to lifting $\epsilon'_0$ to $X^G(S')$, where $L'_0=\uHom_{\OO_{S'_0}}(\epsilon'_0{}^*(\Omega^1_{X'_0/S'_0}),\mathscr I')$. Since $L'_0$ is a $k$-vector space, one has a canonical isomorphism $H^1(G'_0,L'_0)\isomto H^1(G'_0\times_{A'_0}k,L'_0)$. It suffices to show that $H^1(G'_0\times_{A'_0}k,L'_0)=0$. Writing $\overline k$ for an algebraic closure of $k$, one has an isomorphism (IX.3.1 in the affine case, lemma 9.11 in general)
\[
H^1(G'_0\times_{A'_0}k,L'_0)\otimes_k\overline k
\isomto H^1(G'_0\times_{A'_0}\overline k,L'_0\otimes_k\overline k).
\]
One then observes that the representation $L'_0\otimes_k\overline k$ is a direct sum of $(\Omega^1_{X'_0\times_k\overline k})^*$. By hypothesis one has $H^1(G'_0\times_{\overline k}A'_0,(\Omega^1_{X'_0\times_k\overline k})^*)=0$, which implies $H^1(G'_0\times_{A'_0}k,L'_0)=0$.
% typo?: c(epsilon_0) lacks a prime; reversed fiber-product base G'_0 times_{bar k} A'_0 and the displayed differentials retained as printed.
\begin{corollary}{9.8}\label{XII.9.8}
Let $G$ be a flat $S$-group scheme of finite presentation acting on a smooth $S$-scheme $X$ of finite presentation. Suppose that $G$ admits a composition series whose factors are of one of the following types:
\begin{enumerate}[label=(\arabic*)]
\item $S$-abelian scheme (\ie $G$ is smooth over $S$ and its fibers are abelian varieties),
\item $S$-group scheme of multiplicative type,
\item finite étale $S$-group scheme of degree invertible on $S$,
\item reductive $S$-group if $S$ is a $\QQ$-scheme.
\end{enumerate}
Then the functor $X^G$ of fixed points of $X$ under $G$ is representable by a closed subscheme of $X$, of finite presentation, smooth over $S$.
\end{corollary}
Indeed, if $1\to G'\to G\to G''\to1$ is an exact sequence of $S$-groups satisfying the hypotheses of the theorem, the $S$-group $G''$ acts on the functor $X^{G'}$ and the functor $X^G$ is nothing other than that of fixed points of $X^{G'}$ under $G''$. Thus one is reduced to the case of an $S$-group of one of the four types above. By passage to the limit, one can suppose $S$ noetherian. One will verify the cohomological criterion stated above for a fiber $G_s$ over a geometric point $s$ of $S$ and a finite-dimensional linear representation $V$ of $G_s$.

In the case of an abelian scheme, every map from $G_s^n$ to $V$ is constant and consequently the $H^i(G_s,V)$ vanish for $i>0$.

In the case of a group scheme of multiplicative type, $G_s$ is a diagonalizable group. Theorem I.5.3.3 shows that the $H^i(G_s,V)$ vanish for $i>0$.

In the next two cases, the argument is the same, since it rests on the semisimplicity of linear representations of $G_s$. Indeed, in the case of a finite étale $S$-group scheme of degree invertible on $S$, $G_s$ is a finite constant group of degree invertible in the residue field $\kappa(s)$; it is well known that every linear representation of $G_s$ is semisimple.

For the reductive case, the residue field $\kappa(s)$ is supposed (algebraically closed) of characteristic zero. One then knows that every linear representation of $G_s$ is semisimple (see [T-Y], theorem 27.3.3).
% typo: missing "que" after "On sait alors" -> supplied in English.
\begin{remark}{9.9}\label{XII.9.9}
This applies in particular to the case of an action of $G/S$ on a smooth $S$-group scheme $H$. If $G$ is a group of multiplicative type acting by inner transformations on $H$ smooth and affine over $S$, one then recovers corollary XI.5.3 asserting the smoothness of the centralizer of $G$ in $H$.
% typo: "tranformations" -> "transformations".
\end{remark}
\begin{remark}{9.10}\label{XII.9.10}
In the case where $S=\Spec(k)$ and $G$ is a linearly reductive algebraic group defined over the algebraically closed field $k$, this result is due independently to Fogarty ([F], theorem 5.4) and Iversen ([I], proposition 1.3).
\end{remark}

The proof of theorem 9.7 uses the following lemma, well known in the case of an affine group scheme (IX.3.1).
\begin{lemma}{9.11}\label{XII.9.11}
Let $G$ be a group scheme over an affine scheme $S=\Spec(A)$, and $f:S'=\Spec(A')\to S$ a flat morphism. Put $G'=G\times_S S'$. Then for every quasi-coherent $G$-$\OS$-module $M$, one has canonical isomorphisms
\[
H^i(G,M)\otimes_A A'\isomto H^i(G',M\otimes_A A')\qquad(i\ge0).
\]
\end{lemma}
Since chains of degree $n\ge1$ for Hochschild cohomology are determined by their value at the point $(\mathrm{id},\ldots,\mathrm{id})\in G(G)\times\cdots G(G)$, one has an isomorphism of $A$-modules $C^n(G,M)\simeq\Gamma(G^n,M\otimes_A\OO_{G^n})$. The cohomology group $H^n(G,M)$ is the $n$th cohomology group of the complex
\[
\begin{aligned}
L&\to\Gamma(G,M\otimes_A\OO_G)\to\Gamma(G^2,M\otimes_A\OO_{G^2})\\
&\to\Gamma(G^3,M\otimes_A\OO_{G^3})\to\cdots
\end{aligned}
\]
Since $A'$ is flat over $A$, by EGA~IV$_1$, 1.7.21, one has
\[
\Gamma(G^n,M\otimes_A\OO_{G^n})\otimes_A A'
\isomto\Gamma(G'{}^n,(M\otimes_A A')\otimes_{A'}\OO_{G'{}^n}).
\]
This implies the desired assertion on taking cohomology.
% typo?: L, rather than M, as first term of complex retained.

\begin{thebibliography}{RG71}
\item[]\nde{Additional references cited in this Exposé.}
\bibitem[Fo73]{XII.Fo73} J.~Fogarty, \textit{Fixed point schemes}, Amer.~J.~Math. \textbf{95} (1973), 35--51.
\bibitem[RG71]{XII.RG71} M.~Raynaud, L.~Gruson, \textit{Critères de platitude et de projectivité}, Invent.~math. \textbf{13} (1971), 1--89.
\bibitem[Iv72]{XII.Iv72} B.~Iversen, \textit{A fixed point formula for action of tori on algebraic varieties}, Invent.~math. \textbf{16} (1972), 229--236.
\bibitem[Ray72]{XII.Ray72} M.~Raynaud, \textit{Flat modules in algebraic geometry}, Compositio Math. \textbf{24} (1972), 11--31.
\bibitem[TY06]{XII.TY06} P.~Tauvel, R.~W.~T.~Yu, \textit{Lie algebras and algebraic groups}, Springer-Verlag, 2006.
\bibitem[We05]{XII.We05} S.~Wewers, \textit{Formal deformation of curves with group scheme action}, Ann.~Inst.~Fourier (Grenoble) \textbf{55} (2005), 1105--1165.
\end{thebibliography}
